% Evans 10.1: printed pp.543--550; English PDF pp.560--567.
% Chinese section start: printed p.384 / scan p.402.
\section{Introduction, Viscosity Solutions}
\label{sec:10-1}
\subsection{Touching replaces integration by parts}
We study
\begin{equation}\label{eq:10-hj}
 u_t+H(D_xu,x)=0,\qquad u(x,0)=g(x),
\end{equation}
where $H:\R^n\times\R^n\to\R$ is continuous. Characteristics may
intersect, so a globally smooth solution is generally unavailable.
The nonlinear expression $H(Du,x)$ has no general distributional
interpretation for merely continuous $u$. At a local extremum of
$u-\phi$, however, the derivative of a smooth test function can replace
that of $u$. This is a pointwise contact argument, not integration by parts.
Write $\mathrm{BUC}(\R^n\times[0,T])$ for bounded uniformly continuous
functions on the whole closed slab; the uniformity includes spatial infinity.
\Needspace{55mm}
\begin{definition}[Initial-value viscosity convention; source \S10.1.1]
A continuous $u$ is a viscosity subsolution in $\R^n\times(0,T)$ if,
whenever $\phi\in C^\infty$ and $u-\phi$ has a local maximum at $(x,t)$,
\[
 \phi_t(x,t)+H(D\phi(x,t),x)\leq0.
\]
A supersolution uses a local minimum and the reverse inequality.
A solution is both. For the source's initial-value class we additionally
require $u\in\mathrm{BUC}$ on each finite slab and $u(\cdot,0)=g$.
The contact need not be strict and its value need not be zero, since a
constant can be added to $\phi$. All contacts in this definition have $t>0$.
\end{definition}
\Needspace{42mm}
\PDEAnnotation{测试方向与极值}{局部最大点对应测试函数从上接触，要求残差不大于零；局部最小点从下接触，要求残差不小于零。导数放在测试函数上，定义不要求 $Du$ 处处存在。}
\Needspace{55mm}
\begin{lemma}[Persistence of strict contacts; editorial detail for \S10.1.1]
Suppose $u_j\to u$ uniformly on a closed ball $B$ in space--time,
$u-\phi$ has a strict maximum at its center $z_0$, and $B$ lies in the
open PDE domain. Maximizers $z_j$ of $u_j-\phi$ on $B$ lie in its
interior eventually and converge to $z_0$.
\end{lemma}
\begin{expandedproof}[the boundary gap and all cluster points]
Continuity and compactness give a positive gap between the central value
and the maximum over $\partial B$. Uniform convergence preserves half
that gap, so $z_j\notin\partial B$. Any convergent subsequence of $z_j$
has a limit maximizing $u-\phi$ on $B$, by uniform convergence.
Strictness makes this limit $z_0$. Compactness then gives convergence of
the entire sequence. The minimum version follows by changing signs.
For a nonstrict maximum, replace $\phi$ by
$\phi+\delta|z-z_0|^2$: the contact is strict and the first derivatives at
$z_0$ are unchanged. For second-order equations one can instead use
$\delta|z-z_0|^4$ to preserve the Hessian as well.
\end{expandedproof}
\Needspace{55mm}
\begin{theorem}[Vanishing viscosity; source motivation, made conditional]
Let $\varepsilon_j\downarrow0$ and let smooth $u_j$ satisfy
$(u_j)_t+H(Du_j,x)-\varepsilon_j\Delta u_j=0$.
If $u_j\to u$ locally uniformly, then $u$ solves the first-order equation
in the viscosity sense in the open domain. If the convergence includes
$t=0$, the initial values converge to $g$, and $u$ is in the BUC class,
then it is a viscosity solution of the initial-value problem.
\end{theorem}
\begin{expandedproof}[the sign of the positive diffusion]
At a strict upper contact, choose $z_j\to z_0$ by the lemma. There
$D u_j=D\phi$, $(u_j)_t=\phi_t$, and
$D_x^2u_j\leq D_x^2\phi$ as symmetric matrices. Hence
\[
 \phi_t(z_j)+H(D\phi(z_j),x_j)
 =\varepsilon_j\Delta u_j(z_j)
 \leq\varepsilon_j\Delta\phi(z_j).
\]
The right side tends to zero because the test Hessian is bounded near
$z_0$. Continuity of $H$ gives the subsolution inequality. At a lower
contact the Hessian order reverses and gives the supersolution inequality.
Strictification treats arbitrary contacts. Convergence at $t=0$ gives the
trace. Local uniform convergence by itself does not prove global boundedness
or uniform continuity, and this theorem assumes the smooth approximants;
it does not construct them or establish their uniform estimates.
\end{expandedproof}
\Needspace{42mm}
\PDEAnnotation{消黏法的条件式结论}{不必控制 $D^2u_j$ 本身：在接触点只需用 $D^2\phi$ 的单侧界。但要得到存在性，还须另证正则化方程可解及紧性，不能由这个极限论证直接宣布一般 $H$ 的存在性。}

\Needspace{70mm}
\subsection{Smooth tests, differentiability, and consistency}
\Needspace{50mm}
\begin{lemma}[$C^1$ tests suffice; editorial expansion of source Theorem 1]
The first-order definition is unchanged if $C^\infty$ test functions
are replaced by $C^1$ test functions.
\end{lemma}
\begin{expandedproof}[mollify the test, not the unknown]
Only one implication needs proof. Strictify a $C^1$ upper contact,
choose a ball compactly inside the PDE domain, and mollify an extension
of $\phi$ on a slightly larger ball. Then $\phi_j\to\phi$ in $C^1$
on the contact ball. Maxima of $u-\phi_j$ are interior and converge
to the contact point by the boundary-gap argument. Apply the smooth-test
inequality and pass to the limit in the first derivatives and $H$.
For a lower contact repeat with minima. A smooth cutoff extends each
local test to the ambient domain without changing it near the contact.
\end{expandedproof}
\Needspace{50mm}
\begin{lemma}[A differentiable function has a $C^1$ touching function;
source \S10.1.2 lemma]
If continuous $u$ is differentiable at $z_0$, it has a $C^1$ strict
upper test and a $C^1$ strict lower test there, both with derivative $Du(z_0)$.
\end{lemma}
\begin{expandedproof}[a radial majorant of the differentiability error]
Subtract the tangent affine function and translate so $z_0=0$,
$u(0)=0$, $Du(0)=0$. On a small ball write
$|u(z)|\leq |z|\rho(|z|)$, where $\rho$ is continuous,
nondecreasing and $\rho(0)=0$. One choice is the maximum over $|w|\leq r$
of $|u(w)|/|w|$, extending that quotient by zero at $w=0$.
Set for $r=|z|$
\[
 \psi(z)=\int_r^{2r}\rho(s)\,ds+r^2.
\]
Then $\psi(z)\geq r\rho(r)+r^2>u(z)$ for $z\ne0$.
For $z\ne0$ its radial derivative is
$2\rho(2r)-\rho(r)+2r\to0$, while $\psi(z)/|z|\to0$.
Thus $\psi$ extends as a $C^1$ function with $D\psi(0)=0$.
It is a strict upper test. The same construction applied to $-u$
gives a lower test for $u$. Add back the affine function and use a cutoff
if a test on the whole ambient domain is desired.
\end{expandedproof}
\Needspace{55mm}
\begin{theorem}[Consistency; source Theorem 1]
A $C^1$ classical solution is a viscosity solution locally. Conversely,
a viscosity solution differentiable at $(x_0,t_0)$ satisfies
$u_t+H(Du,x_0)=0$ at that point. A locally Lipschitz viscosity solution
therefore solves the equation a.e.
\end{theorem}
\begin{expandedproof}[two contacts recover equality]
For a classical solution, derivatives of $u$ and its test agree at any
interior contact, so both test inequalities follow with equality.
Conversely, apply the preceding lemma in all $n+1$ variables.
The $C^1$ test equivalence gives the PDE residual at most zero from
the upper test, and at least zero from the lower test. Their derivatives
are those of $u$, proving equality. Rademacher's theorem supplies a.e.
differentiability of a locally Lipschitz function. The converse implication
from a.e. solvability to viscosity solvability is false: corners carry
additional information, as the next diagram and Problem 5 demonstrate.
\end{expandedproof}
\PDEBodyDiagram{10.1}{A corner distinguishes a.e. and viscosity solutions.}
{On $(-1,1)$ both $1-|x|$ and $|x|-1$ solve $|u'|=1$ away from zero and vanish at the endpoints. The upper test $1-x^2/2$ is allowed for the downward corner, with slope zero and residual $-1\leq0$. For the upward corner, the lower test $-1+x^2/2$ has the same residual and violates the required supersolution inequality. No smooth lower test exists at the downward corner.}
{ch10-contact}
\Needspace{42mm}
\PDEAnnotation{几乎处处满足还不够}{尖点集合可以只有一个点、测度为零，却决定可接受的解。黏性定义在尖点也检查所有可能接触斜率；只在可微点代入方程会漏掉这个选择机制。}
